\section{Introduction}
\label{sec:intro}

The memory cost of transformer inference grows with context: every decoded
token appends to a key--value cache, and long-horizon reasoning workloads
pay for that growth in high-bandwidth memory. One line of work manages the
growing cache by evicting or compressing entries
\citep{xiao2024streamingllm, zhang2023h2o}. A second line removes the
growth altogether: linear-attention and fast-weight architectures carry a
fixed $d_{\mathrm{state}} \times d_{\mathrm{state}}$ matrix state that is
updated in place as the sequence streams
\citep{katharopoulos2020linear, schlag2021linear, yang2024deltanet,
gu2023mamba}. A fixed-size state is the constant-memory endpoint of the
cache trade-off, and it turns a scheduling question into a provisioning
question: for a state of a given byte size, how many independent bindings
can it be trusted to hold, and how does that number move when the byte
budget changes?

Existing answers do not settle the provisioning question. Classical linear
associative memories have information-theoretic ceilings near
$K \approx d$ for $K$ independent bindings under exact linear reads
\citep{kohonen1972correlation, anderson1972simple}, and recent theory
sharpens such thresholds for linear associative recall
\citep{barnfield2026sharp}. Under argmax decoding the accounting changes
in kind: a rank-one state can recover on the order of $d$ associations
\citep{nichani2025factual}, so any capacity number is meaningful only
relative to its readout. Descriptive studies track the effective rank of
trained linear-attention states at scale
\citep{nazari2026rank, sun2026staterank}. What none of these supply is a
\emph{located} frontier for a \emph{trained} delta-rule memory under an
exact continuous readout: the load at which recall collapses, measured
with a confidence interval, and its scaling in $d_{\mathrm{state}}$.

This paper supplies that frontier on a controlled testbed, with three
results and one provisioning law.

\textbf{First, the transition, located.} At $d_{\mathrm{state}}=64$,
held-out compositional recall does not decline gracefully with load: it
transitions sharply, and a logistic fit over the measured loads places the
midpoint at $x_0 = K/d_{\mathrm{state}} = 0.5455$ (95\% seed-level
bootstrap CI $[0.5385, 0.5513]$, width $0.0127$), % evidence: C2
a roughly $20\times$ tightening over the prior bracket between the two
flanking archived loads. % evidence: C2

\textbf{Second, the ratio is not the law.} The identical
$K/d_{\mathrm{state}}$ window that collapses at $d_{\mathrm{state}}=64$ is
flat at exact ceiling at $d_{\mathrm{state}}=128$: $h4=1.0$ in 12 of 12
cells. % evidence: C3
The frontier at $d_{\mathrm{state}}=80$ sits at $x_0 = 0.6779$ (CI
$[0.6683, 0.6867]$), excluding a pre-registered ratio-invariance band
$[0.4745, 0.6165]$ entirely, % evidence: C10
and $d_{\mathrm{state}}=96$ shows no transition anywhere in a window
reaching $K/d_{\mathrm{state}} = 0.9375$. % evidence: C11
Between the two located frontiers, capacity in raw bindings grows as
$d^{1.97}$ (CI-propagated two-point range $[1.86, 2.09]$): consistent with
proportional to state \emph{bytes} ($d^2$), excluding proportional to
state \emph{width} ($d^1$). % evidence: C15
A load rule calibrated at one scale therefore mis-provisions another;
the safe load ratio itself grows with $d_{\mathrm{state}}$.

\textbf{Third, the leading confound, tested and exonerated.} The collapsing
and non-collapsing scales differ in more than size: with $n=107$ entities,
the $d=64$ anchor table is forced non-orthogonal by the Welch bound while
the $d=128$ table is exactly orthogonal, making table coherence the most
natural single-variable account. Injecting controlled, frozen coherence
into the flat $d=128$ geometry, at and above the $d=64$ table's own
realized training band, leaves recall at exact ceiling in 19 of 19 cells
under both a concentrated (rank-4) and a diffuse (rank-48) injection
structure. % evidence: C6
Scalar table coherence, at matched dose and matched geometry, does not
reproduce the transition.

\paragraph{A binding count as a reasoning-memory proxy.} Multi-step
reasoning traces carry intermediate results as sub-goals or variable
bindings that later steps must retrieve: a chain that derives seven
intermediate quantities before its answer holds seven bindings in
working memory at once, the same simultaneous-binding load this testbed
isolates directly. A fixed-size fast-weight state that keeps
intermediate reasoning state in place, rather than replaying it through
a growing key--value cache, inherits exactly the load limit measured
here: the located frontier converts a state's byte budget into a
concrete bound on how many sub-goals it can hold before recall
degrades. The synthetic grammar does not model any particular reasoning
task, but the binding-count abstraction it isolates is the quantity a
constant-memory reasoning architecture would need to provision.

Every wave reported here was pre-registered before its cells ran, with
outcome semantics fixed in advance, including the rule that a degenerate
sigmoid fit on flat data must be disclosed as such and never quoted as a
located transition; instrument saturation is checked at every wave rather
than assumed. The scope is disclosed up front: a single-head, single-layer
DeltaNet, a 107-entity synthetic grammar, sub-million-parameter models,
and $d_{\mathrm{state}} \in \{64, 80, 96, 128\}$. No claim is made about
production-scale models; Section~\ref{sec:limitations} states the
boundaries.
